% ===================================================================
%  TABEL Nr. 9 — Mägi. Mets.
%  Traduction 2026 d'apres texto/io, controlee sur texto/fr et
%  texto/fi. Consigne : voir texto/et/10-tabel-01.tex.
%
%  CETTE PAGE EST L'EXACT NEGATIF DU TABLEAU 9 ROMANCHE.
%
%  Le romanche y donnait a l'Europe « piz », « lavina »,
%  « muntanella » : il avait la montagne, et c'est lui qui avait
%  fourni les mots. L'ESTONIE N'A PAS DE MONTAGNE. Son point le plus
%  haut fait trois cent dix-huit metres et s'appelle
%  « Suur Munamägi », la grande colline-en-oeuf. « mägi » designe la
%  colline et la montagne du meme mot ; le glacier, « liustik », est
%  une derivation faite au XXe siecle sur « liuguma » glisser ; et
%  l'avalanche est « laviin », emprunt. Sur toute la premiere scene,
%  l'estonien n'a rien a lui.
%
%  MAIS LA SECONDE SCENE EST UNE FORET, ET LA C'EST L'INVERSE. La
%  moitie du pays est boisee, et l'estonien nomme un bois par son
%  arbre au moyen D'UN SEUL SUFFIXE, « -ik » :
%    « männik » la pineraie, de « mänd » le pin ;
%    « kuusik » la sapiniere, de « kuusk » ;
%    « kaasik » la boulaie, de « kask » ;
%    « tammik » la chenaie, de « tamm » ;
%    « haavik » la tremblaie, de « haab » ;
%    « lepik »  l'aunaie, de « lepp ».
%  L'allemand doit composer — Kiefernwald, Birkenwald —, le francais
%  doit phraser — bois de pins —, et l'ido de cette page ecrit « bosko
%  de abieti », trois mots.
%
%  ET C'EST LA SECONDE FORME DE LA REGLE DU TABLEAU 7 ROMANCHE. On y
%  avait vu qu'une RACINE PLEINE marque ce qu'on doit distinguer en
%  agissant — « bov », « tor », « vatga », « vadè ». Mais une racine
%  ne vaut que pour un ensemble ferme : il y a quatre etats d'un
%  bovin et l'on peut les apprendre. Les especes d'arbres, elles,
%  n'ont pas de fin. POUR UNE DISTINCTION OUVERTE, IL NE FAUT PAS UNE
%  RACINE, IL FAUT UN SUFFIXE — quelque chose qui se pose sur
%  n'importe quel nom d'arbre et qui donne le bois sans qu'on ait
%  rien a retenir.
%
%  Racine pour ce qui se compte, suffixe pour ce qui ne se compte
%  pas : c'est la meme economie, et l'homme qui marche en foret tous
%  les jours ne veut pas d'une periphrase par essence.
% ===================================================================

%%K t09-apar-1 apar
\VUsaut{4.0mm}
\VUtitre{15.0pt}{60}{TABEL Nr. 9}
\VUsaut{3.0mm}

%%K t09-tit-1 sub
\VUcentre{12.0pt}{0}{\textit{Esimene stseen.}}
\VUsaut{3.0mm}

%%K t09-tit-2 sub
\VUpk{11.4pt}{40}{Mägi. --- Kuurortlinn.}
\VUsaut{3.0mm}

%%K t09-c1-01-1 p
1. --- Kaheksa päeva olen ma Pratzis. Ma ei suuda väljendada kogu imetlust, mida ma
tundsin, kui ma olin Püreneede imepärase \VUgras{mäeaheliku} \textsuperscript{(1)} ees,
mille \VUgras{hari} \textsuperscript{(2)} paistab sinises taevas nagu hiiglaslik
pärjand.

%%K t09-c1-02-1 p
2. --- Ma ei kujutanud ette, et Püreneede \VUgras{tipud} \textsuperscript{(3)} ulatuvad
nii suure \VUgras{kõrguseni}, ja ma jäin jahmatusse imeilusate \VUgras{liustike}
\textsuperscript{(5)} ja lumiste \VUgras{mäetippude} \textsuperscript{(4)} ees, millelt
veerevad nii hirmuäratavad \VUgras{laviinid}.

%%K t09-tit-3 sub
\VUsaut{3.0mm}
\VUpk{10.2pt}{40}{Mäemaastik.}
\VUsaut{3.0mm}

%%K t09-c1-03-1 p
3. --- Juba oma saabumisele järgnenud päeval läksin ma vana kustunud \VUgras{vulkaani}
\textsuperscript{(6)} juurde, mis on Pratzi lähedal. Kraatri paljastel ja lagedatel
äärtel on veel hästi näha jäljed \VUgras{mäeküljel} \textsuperscript{(7)} mitu sajandit
tagasi voolanud laava teekonnast. Mul õnnestus ühelt \VUgras{nõlvalt} leida mõned tükid
sellest laavast.

%%K t09-c1-04-1 p
4. --- Pratz asub piiril, ja \VUgras{kindlus} \textsuperscript{(8)}, mis kaitseb
Prantsusmaad Hispaaniast eraldavat \VUgras{kuru} \textsuperscript{(27)}, meenutab täiesti
neid vanu Reini kalda losse, mis näivad oma kalju otsast \VUgras{saaki} varitsevat.

%%K t09-c1-05-1 p
5. --- Tohutu \VUgras{kotkas} \textsuperscript{(9)}, kellel on pesa mitte kaugel
\VUgras{kindlusest} \textsuperscript{(8)}, köidab sageli mu tähelepanu.

See \VUgras{röövlind}, kelle \VUgras{nokk} on tugev, toob sageli oma poegadele talle või
kitsetalle, keda ta hoiab oma \VUgras{küüntega}. Nii suur on tema \VUgras{tiibade}
ulatus, et ta suudab kaua liuelda oma raske koormaga.

%%K t09-tit-4 sub
\VUsaut{3.0mm}
\VUpk{10.2pt}{40}{Matkaja.}
\VUsaut{3.0mm}

%%K t09-c1-06-1 p
6. --- Seal on kõige meeldivam jalutuskäik retk «Cau» koski \textsuperscript{(10)}
juurde. \VUgras{Veelangus} kukub keeristes ja kaob siis ilusa \VUgras{ojana}
\VUgras{kuusikusse} \textsuperscript{(11)}, mis asub linnast läänes. Me tõusime läbi selle
kuusiku \VUgras{meteoroloogiajaamani} \textsuperscript{(12)} ja
\VUgras{vaateplatvormi} otsast nägime kogu piirkonna \VUgras{panoraami}.
\VUgras{Maastik} on imepärane, ja \VUgras{loodus} näib maale pillavat kõiki aardeid, mis
tal on.

%%K t09-c1-07-1 p
7. --- Alla minekuks kasutasime me \VUgras{köisraudteed} \textsuperscript{(13)} ja
peatusime väikeses \VUgras{jaamas} vana \VUgras{feodaallossi} \VUgras{varemete}
\textsuperscript{(14)} juures, kus kunagi elasid Pratzi isandad.

%%K t09-c1-08-1 p
8. --- Vaene loss! Mida ta küll mõtleb, nähes all kenat \VUgras{kuurortlinna}
\textsuperscript{(15)}, mis on nüüd moodne \VUgras{termaalkuurort}?

%%K t09-c1-08-2 p
\VUgras{Kliima} on nii meeldiv, \VUgras{mineraalvesi} nii mõjus, et
\VUgras{sanatooriumis} \textsuperscript{(17)} kasutatav \VUgras{ravi} on tõeline
lõbu.

%%K t09-tit-5 sub
\VUsaut{3.0mm}
\VUpk{10.2pt}{40}{Meeldiv termaallinn.}
\VUsaut{3.0mm}

%%K t09-c1-09-1 p
9. --- Pealegi on tehtud kõik, et viibimine meeldivaks teha. Suur \VUgras{viadukt}
\textsuperscript{(16)} ehitati raudtee võimaldamiseks; \VUgras{termaalasutuse}
\VUgras{park} \textsuperscript{(18)}, kus \VUgras{antakse kontserte}, on võluvalt
seatud. Mitu graatsilist «\VUgras{villat}» \textsuperscript{(19)} ehitati kasiino
\textsuperscript{(21)} ümber, ja väga hästi hooldatud \VUgras{maantee}
\textsuperscript{(22)} hõlbustab suuresti läbikäimist.

%%K t09-c1-10-1 p
10. --- Lihtne \VUgras{nõlvak} \textsuperscript{(23)} eraldab \VUgras{teed}
\textsuperscript{(20)} päris \VUgras{orust} \textsuperscript{(25)}, mille põhjas on
graatsiline väike \VUgras{järveke} \textsuperscript{(26)}, mis toidab selget \VUgras{oja}
\textsuperscript{(24)}.

%%K t09-c1-11-1 p
11. --- Oru teisel pool käib \VUgras{rauakaevanduse} \textsuperscript{(28)} töö. Mitu
korda käisin ma vaatamas \VUgras{kaevureid}, kes laskuvad \VUgras{šahti}, ja ma isegi
sisenesin \VUgras{käiku}, kus nad veedavad oma päeva, kaevandades \VUgras{maaki}, mis
sisaldab \VUgras{metalli}.

%%K t09-c1-12-1 p
12. --- Kaevanduse lähedale ehitati tähtis \VUgras{sulatustehas}, et kohe ära kasutada
sealt saadavat rikkust. \VUgras{Kõrgahi} \textsuperscript{(29)}, mida soojendab siin
rohkesti leiduv \VUgras{kivisüsi}, võimaldab kiiresti ja odavalt saada \VUgras{malmi}.

%%K t09-c1-13-1 p
13. --- Mitte kaugel kaevandusest kaevatakse \VUgras{kivimurdu} \textsuperscript{(30)}.
Ei \VUgras{graniiti}, ei \VUgras{porfüüri}, ei isegi \VUgras{liivakivi} ei raiu vana
\VUgras{kivitööline} oma \VUgras{kirkaga}, vaid lihtsat \VUgras{tahutud kivi} majade
ehitamiseks.

%%K t09-c1-14-1 p
14. --- Üks selle maa ilusamaid ehteid on \VUgras{kuused} \textsuperscript{(31)}. Kõikjal
\VUgras{kuristiku} \textsuperscript{(32)} põhjas, \VUgras{mägijõe} \textsuperscript{(33)},
\VUgras{sügaviku} \textsuperscript{(34)} või järsaku kaldal annavad nende peened okkad
rohelise varjundi.

%%K t09-tit-6 sub
\VUsaut{3.0mm}
\VUpk{10.2pt}{40}{Jalgsi käimine.}
\VUsaut{3.0mm}

%%K t09-c1-15-1 p
15. --- Ma kasutan oma vaheaega \VUgras{pikkadeks} jalutuskäikudeks. \VUgras{Teed}
\textsuperscript{(35)}, mis looklevad mäel, võimaldavad läbida, distantsile mõtlemata,
mitu \VUgras{kilomeetrit} ja sageli mitu \VUgras{kümmet kilomeetrit}.

%%K t09-c1-15-2 p
\VUgras{Piirikivid} \textsuperscript{(36)} ja \VUgras{viidapostid} \textsuperscript{(37)}
tähistavad teid, millel kohtab sageli noort \VUgras{mägilast} \textsuperscript{(38)}
\VUgras{kandekotiga} \textsuperscript{(40)} küljel, kes kannab \VUgras{marmotit}
\textsuperscript{(39)}, või mõnda \VUgras{mustlast} \textsuperscript{(41)}, kelle
\VUgras{karusnahkne kuub} \textsuperscript{(42)} on kummalises vastuolus vana rebenenud
\VUgras{kapuutsiga} \textsuperscript{(43)}. Ta talutab ketiga taltsutatud \VUgras{karu}
\textsuperscript{(44)}.

%%K t09-tit-7 sub
\VUpk{10.2pt}{40}{Mägedele ronimine.}
\VUsaut{3.0mm}

%%K t09-c1-16-1 p
16. --- Peaaegu iga päev käin ma \VUgras{teejuhiga} \textsuperscript{(48)}
\VUgras{ronimas}, ja ma olen väga uhke, kui ma, \VUgras{seljakott}
\textsuperscript{(47)} seljas ja «\VUgras{alpikepp}» käes, nagu tõeline \VUgras{turist}
\textsuperscript{(46)}, ründan \VUgras{kaljusid} \textsuperscript{(45)}.

%%K t09-c1-17-1 p
17. --- Mäe otsast paistavad tasandiku elanikud mitte suuremad kui sipelgad;
\VUgras{lambakari} \textsuperscript{(50)}, mis karjatab \VUgras{karjamaal}
\textsuperscript{(49)}, näib laste mänguasjana. \VUgras{Mänd} \textsuperscript{(51)} või ka
\VUgras{puumajake} \textsuperscript{(52)} paistab vaevalt täpina rohelusel.

%%K t09-c1-18-1 p
18. --- Nende retkede ajal kohtame me sageli \VUgras{rajal} \textsuperscript{(53)}
\VUgras{mägikitsekütti} \textsuperscript{(54)}, kes kannab õlal looma, kelle ta äsja
tappis.

%%K t09-c1-19-1 p
19. --- Ma panin oma herbaariumi väga tähelepanuväärse \VUgras{sõnajala}
\textsuperscript{(55)} oksa ja ilusa roosa \VUgras{kanarbiku} \textsuperscript{(57)} oksakese,
mille ma korjasin väikese \VUgras{koopa} \textsuperscript{(56)} suudmes oma viimasel
jalutuskäigul.

%%K t09-tit-8 sub
\VUsaut{3.0mm}
\VUcentre{12.0pt}{0}{\textit{Teine stseen.}}
\VUsaut{3.0mm}

%%K t09-tit-9 sub
\VUpk{11.4pt}{40}{Mets. --- Jaht.}
\VUsaut{3.0mm}

%%K t09-c2-01-1 p
1. --- Eile veetsin ma metsas imelise päeva. Usinus, mis valitseb loomade ja
\VUgras{putukate} \textsuperscript{(5)} seas, kes seda asustavad, huvitas mind väga.

%%K t09-c2-01-2 p
\VUgras{Ämblik} \textsuperscript{(1)}, kes koob oma \VUgras{võrku} \textsuperscript{(2)}
kahe oksa vahele, et püüda mööduvat \VUgras{kärbest} \textsuperscript{(3)}, \VUgras{tigu}
\textsuperscript{(4)}, kes vaevaliselt kannab oma koda, põdrapõrnikas, kes oma saaki
uurib, usin sipelgas, kes on väga innukas \VUgras{sipelgapesa} \textsuperscript{(6)}
juures, kõik need väikesed käisid, tulid, töötasid erakordse usinusega.

%%K t09-c2-02-1 p
2. --- \VUgras{Madu} \textsuperscript{(7)}, keda ma esimesel pilgul pidasin
\VUgras{rästikuks}, kuid kes ei olnud midagi muud kui lihtne \VUgras{nastik}, oli
peitunud puutüve alla ja pistis aeg-ajalt välja oma peenikest peakest, mis näis alati
hammustamiseks valmis. Minu ees roomas raskelt jäme \VUgras{nälkjas}
\textsuperscript{(8)}, olles ära söönud ühe maitsvatest \VUgras{maasikatest}
\textsuperscript{(11)}, mida seal rohkesti kasvab.

%%K t09-c2-03-1 p
3. --- Maapind oli vaibaks kaetud \VUgras{metsalilledega}: \VUgras{nurmenukud}
\textsuperscript{(9)}, \VUgras{igihaljad} \textsuperscript{(10)} ja palju muid taimi, mida
ma ei tundnud, õitsesid \VUgras{rohutuustide} \textsuperscript{(12)} vahel.

%%K t09-c2-04-1 p
4. --- Selles piirkonnas on \VUgras{trühvel} peaaegu sama tavaline kui \VUgras{seen}
\textsuperscript{(14)}, ja metsavahile kuuluv \VUgras{põrsas} \textsuperscript{(13)} uuris
maiuslikult, kas tal õnnestub mõnda neist leida.

%%K t09-tit-10 sub
\VUsaut{3.0mm}
\VUpk{10.2pt}{40}{Salaküttimine.}
\VUsaut{3.0mm}

%%K t09-c2-05-1 p
5. --- Ma viibisin stseeni \VUgras{lõpus}, mis mind väga lõbustas. Oma
\VUgras{lühijalgse koera} \textsuperscript{(15)} haistmise abil oli \VUgras{metsavaht}
\textsuperscript{(16)} äsja üllatanud \VUgras{salaküti} \textsuperscript{(22)} hetkel, mil
see valmistus ära viima \VUgras{ulukit} \textsuperscript{(17)}, mille ta oli tapnud.
Eelistades oma saaki maha jätta kui end vahistada lasta, oli salakütt põgenenud, ja
\VUgras{jänes} \textsuperscript{(18)}, \VUgras{emahirv} \textsuperscript{(19)},
\VUgras{metskits} \textsuperscript{(20)} ja \VUgras{metssiga} \textsuperscript{(21)}
lamasid rohul, rõõmustades metsavahti.

%%K t09-c2-06-1 p
6. --- Metsa puude mitmekesisus on hämmastav. \VUgras{Pöögid} \textsuperscript{(23)} ja
\VUgras{kased} \textsuperscript{(26)}, \VUgras{männid}, mis annavad \VUgras{vaiku}
\textsuperscript{(27)}, \VUgras{tammed} \textsuperscript{(29)} ja veel palju teisi segavad
oma oksi.

%%K t09-c2-06-2 p
\VUgras{Luuderohi} \textsuperscript{(24)}, mis klammerdub kõrgete puude tüve külge, annab
\VUgras{metsale} \textsuperscript{(25)} sädeleva rohelise, mis ühineb meeldivalt
\VUgras{sambla} \textsuperscript{(31)} heledama ja kahvatuma rohelise tooniga, milles
\VUgras{roheline rähn} \textsuperscript{(30)} putukaid otsib.

%%K t09-tit-11 sub
\VUsaut{3.0mm}
\VUpk{10.2pt}{40}{Metsamaastik.}
\VUsaut{3.0mm}

%%K t09-c2-07-1 p
7. --- Vanad tammed võlusid mind. Nende jämedad väändunud \VUgras{juured}
\textsuperscript{(32)}, poolenisti mullast välja tulnud, annavad neile uhke jõu ja
tugevuse ilme, ja ma küsin endalt, kuidas võib \VUgras{omanik} \textsuperscript{(35)}
otsustada lasta need maha võtta ja anda need \VUgras{saele} \textsuperscript{(34)}, mis
kuulub \VUgras{raiujale} \textsuperscript{(33)}. Vaesed \VUgras{tüved}
\textsuperscript{(36)}, oma \VUgras{koorest} \textsuperscript{(37)} paljaks tehtud või
\VUgras{kirvega} \textsuperscript{(38)} lõhestatud, mida hoiab \VUgras{päeviline}
\textsuperscript{(39)}, teevad mulle kurvaks.

%%K t09-c2-08-1 p
8. --- Lähedal on vana \VUgras{söepõletaja} \textsuperscript{(41)} valmistanud oma
\VUgras{labidaga} söe\VUgras{kuhja} \textsuperscript{(42)}, ja vanaeit tõi
\VUgras{kubu} \textsuperscript{(40)} \VUgras{kuivanud puitu} mahavõetud puude okstest.

%%K t09-tit-12 sub
\VUsaut{3.0mm}
\VUpk{10.2pt}{40}{Jaht.}
\VUsaut{3.0mm}

%%K t09-c2-09-1 p
9. --- Ma tahtsin minna mõneks hetkeks puhkama \VUgras{metsavahi majja}
\textsuperscript{(46)}, kuid kui ma jõudsin \VUgras{järvekese} \textsuperscript{(43)}
juurde, millel ujub ilus \VUgras{luik} \textsuperscript{(45)}, märkasin ma, et hiljutine
\VUgras{üleujutus} \textsuperscript{(44)} oli tee tõeliseks \VUgras{sooks} muutnud. Ma
pidin ümber pöörama, ja möödudes \VUgras{seedrist} \textsuperscript{(49)},
\VUgras{paplitest} \textsuperscript{(48)} ja \VUgras{jalakast} \textsuperscript{(47)}, mis
on metsa serval, nägin ma \VUgras{võsast} \textsuperscript{(54)} välja tulemas imeilusat
\VUgras{hirve} \textsuperscript{(50)}, keda taga ajas visa \VUgras{jahikoerte kari}
\textsuperscript{(51)}, keda ergutas ka \VUgras{sarv} \textsuperscript{(53)}, mida puhusid
\VUgras{ratsakütid} \textsuperscript{(52)}. Vaene loom oli täiesti kurnatud. Ta kukkus
surres mõne sammu kaugusele minust.

%%K t09-c2-10-1 p
10. --- Metsavahi poeg, keda \VUgras{ajujahi} lärm oli ligi tõmmanud, tuli majast välja
ja me läksime kahekesi tagasi metsa. Ta oli näinud \VUgras{harakat}
\textsuperscript{(56)} vana tamme oksal. Ta arvas, et selle pesa on arvatavasti peidus
\VUgras{lehestikus} \textsuperscript{(58)}, ja ta tahtis pesa kätte saada.

%%K t09-c2-10-2 p
Väle nagu \VUgras{orav} \textsuperscript{(61)}, ronis ta \VUgras{oksalt}
\textsuperscript{(55)} oksale, kuid ta ei leidnud midagi ja tal õnnestus ainult lendu
ajada vana \VUgras{vares} \textsuperscript{(60)}, kes istus \VUgras{harul}
\textsuperscript{(57)}.

\VUsaut{4.0mm}
\VUfilet{22mm}
